%% Preambulo LaTeX: Define classes e características do documento




%% Tipografia
%% Abra este arquivo e selecione uma das opções de fonte nele. A padrão é Times.
\input{configuracoes/tipografia}



%% Pacotes usados pelo documento (se não entender não mexa, hehehe)
\usepackage{courier}                    % Permite a utilização da fonte Courier (para códigos-fonte)
\usepackage[T1]{fontenc}				% Seleção de códigos de fonte.
\usepackage[utf8]{inputenc}				% Codificação do documento (conversão automática dos acentos)
\usepackage{indentfirst}				% Indenta o primeiro parágrafo de cada seção.
\usepackage{nomencl} 					% Usado pela Lista de símbolos
\usepackage{color}						% Controle das cores
\usepackage{graphicx}					% Inclusão de gráficos
\usepackage{float}						% Melhorias para posicionamento de gráficos e tabelas
\usepackage{microtype} 					% Melhorias na justificação
\usepackage{lastpage}   		        % Dá acesso ao número da última página do documento
\usepackage{booktabs}					% Comandos para tabelas
\usepackage{multirow, array}			% Múltiplas linhas e colunas em tabelas
\usepackage{titlesec}                   % Permite criar múltiplas sub seções
\usepackage[table,xcdraw]{xcolor}       % Cores para algumas tabelas especiais
\usepackage[brazilian,hyperpageref]{backref}	 % Inclui nas Referências as páginas onde há as citações
\usepackage{simplecd}                   % Pacote para gerar capa do CD
\usepackage[final]{pdfpages}            % Pacote para incluir um PDF dentro de outro (ficha catalográfica)
%\usepackage{abntex2}



%% Adiciona as alterações do ABNTeX-IFPI
\usepackage{abntex-ifpi/abntex-ifpi}

% Modificações do ABNTeX para o IFPI
\usepackage{abntex-ifpi/tikz-uml}	    % Pacote Tikz UML para criar UML no LaTeX





%% Metadados
\input{configuracoes/pdf}



%% Metadados
%% %%%%%%%%%%%%%%%%%%%%%%%%%%%%%%%%%%%%%%%%%%%%%%%% %%
%% Metadados do trabalho
%% AVISO: Todos esses dados serão automaticamente convertidos para caixa alta onde necessário
%% %%%%%%%%%%%%%%%%%%%%%%%%%%%%%%%%%%%%%%%%%%%%%%%% %%

%% Título
\titulo{CalPEC: sistema eletrônico para cálculo do processo de execução criminal}

%% Autor
\autor{Kaylanne Mendes Dos Santos}

%% Nome do Curso (usado para a Capa do CD)
\nomedocurso{Análise e Desenvolvimento de Sistemas}

%% Local de publicação \local{Teresina, Piauí}

%% Preâmbulo do trabalho
\preambulo{Trabalho de Conclusão de Curso (monografia) apresentado como exigência parcial para obtenção do diploma do Curso de Análise e Desenvolvimento de Sistemas do Instituto Federal de Educação, Ciência e Tecnologia do Piauí, Campus Teresina Central.}

%% Orientador
%% "M\textsuperscript{e}." = Abreviação oficial para "Mestre"
\orientador{Prof. Me. Fernando Castelo Branco Gonçalves Santana}

%% Tipo de Trabalho
%% - Monografia
%% - Tese (Mestrado)
%% - Tese (Doutorado)
%% - Relatório técnico
\tipotrabalho{Monografia}

%% Data do Trabalho \data{2026}

%% Nome da Instituição (para a capa)
\instituicao{INSTITUTO FEDERAL DE EDUCAÇÃO, CIÊNCIA E TECNOLOGIA DO PIAUÍ
\\
CAMPUS TERESINA CENTRAL
\\
TECNOLOGIA EM ANÁLISE E DESENVOLVIMENTO DE SISTEMAS}

%% Primeiro membro da banca examinadora
%%\membroum{Prof. M\textsuperscript{e}. Nome Primeiro Membro da Banca}
\membroum{ Prof. Dr. Ricardo Martins Ramos }

%% Segundo membro da banca examinadora
\membrodois{Prof. Dr. Nadia Mendes dos Santos}

%% Terceiro membro da banca examinadora
%\membrotres{Prof. Thiago Alves Elias da Silva}

%% Data da apresentação do trabalho
%% Se não souber a data da apresentação, utilize \underline{\hspace{3.5cm}}
%% Isso cria um sublinhado de 3.5cm, onde você pode escrever a data depois!
%\dataapresentacao{02 de Outubro de 2023}
\dataapresentacao{18/06/2026}




%% Configuração do "Citado nas páginas"
\input{configuracoes/citacoes}



%% Cores
\input{configuracoes/cores}



%% Espaçamentos
%% Espaçamentos
%% O tamanho do parágrafo é dado por:
\setlength{\parindent}{1.5cm}

%% Espaçamento entre um parágrafo e outro:
%% O abntex diz: "tente também \onelineskip"
\setlength{\parskip}{0cm}

%% Espaçamento entre títulos de seção e o texto/parágrafo seguinte
%% (Formato: {espaço à esquerda}{espaço ANTES do título}{espaço DEPOIS do título, antes do parágrafo})
\titlespacing*{\section}{0pt}{2\baselineskip}{1\baselineskip}
\titlespacing*{\subsection}{0pt}{2\baselineskip}{1\baselineskip}
\titlespacing*{\subsubsection}{0pt}{2\baselineskip}{1\baselineskip}
